\documentclass[11pt]{article}
\usepackage{mathblatt}


\begin{document}
\pfheftstil
\pfheftkopf{Daten}{P10}
\pfrueckblick
\begin{pfaufg}{}{1.}{}
Ordne der Größe nach: 52,2; 19,2; 85,1; 35,3.
\fusshilfe{\mbox{1:~19,2 < 35,3 < 52,2 < 85,1}}
\end{pfaufg}
\begin{pfaufg}{}{2.}{}
Berechne: 1400 : 7; 528 : 11; 16 : 4.
\rechenplatz{1}
\fusshilfe{\mbox{2:~200; 48; 4}}
\end{pfaufg}
\begin{pfaufg}{}{3.}{}
Ergänze die Tabelle Anteil | Winkel.
\par\smallskip\noindent \begingroup\renewcommand{\arraystretch}{1.3}\begin{tabular}{|c|c|}\hline \textbf{Anteil} & \textbf{Winkel} \\ \hline 360° & \rule{0pt}{3ex}\hspace{1.6cm} \\ \hline 180° & \rule{0pt}{3ex}\hspace{1.6cm} \\ \hline 90° & \rule{0pt}{3ex}\hspace{1.6cm} \\ \hline 36° & \rule{0pt}{3ex}\hspace{1.6cm} \\ \hline 3,6° & \rule{0pt}{3ex}\hspace{1.6cm} \\ \hline\end{tabular}\endgroup\par
\fusshilfe{\mbox{3:~360°; 180°; 90°; 36°; 3,6°}}
\end{pfaufg}
\begin{pfaufg}{}{4.}{}
Eine Achse reicht mit 8 Kästchen bis 160. Wie viel ist ein Kästchen? Wie viele Kästchen hoch ist eine Säule für 90?
\rechenplatz{1}
\fusshilfe{\mbox{4:~20 je Kästchen; 4,5 Kästchen}}
\end{pfaufg}
\begin{pfaufg}{}{5.}{}
Ergänze die Tabelle Wort | Prozent für: die Hälfte, ein Viertel, drei Viertel, jede 15.
\par\smallskip\noindent \begingroup\renewcommand{\arraystretch}{1.3}\begin{tabular}{|c|c|}\hline \textbf{Wort} & \textbf{Prozent} \\ \hline 50 \% & \rule{0pt}{3ex}\hspace{1.6cm} \\ \hline 25 \% & \rule{0pt}{3ex}\hspace{1.6cm} \\ \hline 75 \% & \rule{0pt}{3ex}\hspace{1.6cm} \\ \hline $\tfrac{1}{15}$ & \rule{0pt}{3ex}\hspace{1.6cm} \\ \hline\end{tabular}\endgroup\par
\fusshilfe{\mbox{5:~50 \%; 25 \%; 75 \%; $\tfrac{1}{15}$ = $\tfrac{20}{3}$ \% $\approx$ 6,7 \%}}
\end{pfaufg}
\pfstufekopf[sminimummaximumspannw]{Minimum, Maximum, Spannweite}{in 5 der letzten 5 Prüfungen}
\pfgruppe{}
\begin{pfaufg}{}{6.}{}
Jemand fragt: „Welche Schuhgröße kommt am häufigsten vor?“ Kreuze an, welche Kenngröße gefragt ist.
\pfkreuzzeile{Spannweite}\pfkreuzzeile{Modalwert}\pfkreuzzeile{Median}\pfkreuzzeile{Mittelwert}
\fusshilfe{\mbox{6:~Modalwert}}
\end{pfaufg}
\pfgruppe{}
\begin{pfaufg}{P10 ’14}{7.}{}
Eine Tankstelle hat im Sommer 2012 zehn Wochen lang einmal pro Woche die Preise für die Kraftstoffe E 10 und Diesel notiert (in Cent pro Liter). Lies aus der Tabelle den niedrigsten und den höchsten Preis für E 10 ab (Minimum und Maximum).
\par\smallskip\noindent \sachtabelle{lcc}{Datum & E 10 & Diesel}{03.07. & 154,2 & 139,6\\ 10.07. & 157,4 & 142,5\\ 17.07. & 158,8 & 144,2\\ 24.07. & 159,7 & 145,4\\ 31.07. & 159,1 & 146,2\\ 07.08. & 161,0 & 147,1\\ 14.08. & 164,8 & 150,7\\ 21.08. & 169,2 & 154,0\\ 28.08. & 168,1 & 153,1\\ 04.09. & 167,0 & 152,0}\par
\fusshilfe{\mbox{7:~Minimum 154,2 ct; Maximum 169,2 ct}}
\end{pfaufg}
\begin{pfaufg}{P10 ’14}{8.}{}
Eine Tankstelle hat im Sommer 2012 zehn Wochen lang einmal pro Woche die Preise für die Kraftstoffe E 10 und Diesel notiert (in Cent pro Liter). Berechne die Spannweite der Dieselpreise in diesen zehn Wochen.
\par\smallskip\noindent \sachtabelle{lcc}{Datum & E 10 & Diesel}{03.07. & 154,2 & 139,6\\ 10.07. & 157,4 & 142,5\\ 17.07. & 158,8 & 144,2\\ 24.07. & 159,7 & 145,4\\ 31.07. & 159,1 & 146,2\\ 07.08. & 161,0 & 147,1\\ 14.08. & 164,8 & 150,7\\ 21.08. & 169,2 & 154,0\\ 28.08. & 168,1 & 153,1\\ 04.09. & 167,0 & 152,0}\par
\fusshilfe{\mbox{8:~14,4 ct}}
\end{pfaufg}
\begin{pfaufg}{NRW ’23}{9.}{}
Amelie spielt Handball. In fünf Spielen warf sie 12, 13, 9, 15 und 6 Tore. Gib die Spannweite und den Median an.
\par\smallskip\noindent \sachtabelle{l}{Spiel 1 bis 5: 12}{13\\ 9\\ 15\\ 6 Tore}\par
\fusshilfe{\mbox{9:~Spannweite 15 $-$ 6 = 9; Median 12}}
\end{pfaufg}
\begin{pfaufg}{P10 ’22}{10.}{}
Von 2017 bis 2021 wurden jedes Jahr 1 200 Jugendliche gefragt, wie oft sie Bücher lesen. Das Säulendiagramm zeigt, wie viele von ihnen mehrmals pro Woche ein Buch lesen. Gib den kleinsten und den größten Wert im Diagramm an (Minimum und Maximum). Bestimme den Durchschnitt (das arithmetische Mittel) der fünf Werte.
\par\smallskip\noindent \begin{tikzpicture}\begin{axis}[ybar,width=8.5cm,height=4.6cm,ymin=390,ymax=508.75,symbolic x coords={2017,2018,2019,2020,2021},xtick=data,enlarge x limits=0.08,bar width=6mm,ylabel={Anzahl},ylabel style={font=\small},tick label style={font=\small},nodes near coords,every node near coord/.append style={font=\scriptsize,/pgf/number format/.cd,use comma,1000 sep={}},ymajorgrids,grid style={mbgitter}]\addplot[fill=mbkasten,draw=black] coordinates {(2017,468.0) (2018,432.0) (2019,485.0) (2020,470.0) (2021,400.0)};\end{axis}\end{tikzpicture}\par
\rechenplatz{2}
\fusshilfe{\mbox{10:~Minimum 400 (2021), Maximum 485 (2019); Mittel 451}}
\end{pfaufg}
\pfgruppe{}
\begin{pfaufg}{P10 ’19}{11.}{}
Gib die Spannweite dieser Daten an: 9,7 m; 9,9 m; 9,6 m; 9,5 m; 9,8 m; 9,7 m
\rechenplatz{1}
\fusshilfe{\mbox{11:~0,4 m}}
\end{pfaufg}
\pfgruppe{}
\begin{pfaufg}{P10 ’24}{12.}{}
Das Säulendiagramm zeigt, wie viel Niederschlag 2021 in Brandenburg in jedem Monat gefallen ist (in mm). Ermittle die Spannweite der monatlichen Niederschläge im Jahr 2021 in mm.
\par\smallskip\noindent \begin{tikzpicture}\begin{axis}[ymin=0,ymax=100,ytick distance=20,tick label style={font=\small},xtick pos=bottom,ytick pos=left,yticklabel style={/pgf/number format/use comma,/pgf/number format/1000 sep={\,},/pgf/number format/fixed,/pgf/number format/precision=1},ymajorgrids,grid style={mbgitter},xtick={0,1,2,3,4,5,6,7,8,9,10,11},xticklabels={{Jan},{Feb},{Mär},{Apr},{Mai},{Jun},{Jul},{Aug},{Sep},{Okt},{Nov},{Dez}},xmin=-0.5,xmax=11.5,bar width=5mm,nodes near coords,every node near coord/.append style={font=\scriptsize,/pgf/number format/use comma,/pgf/number format/1000 sep={\,},/pgf/number format/fixed,/pgf/number format/precision=1},ybar,width=13.5cm,height=5.2cm,ylabel={Niederschlag in mm},ylabel style={font=\small}]\addplot[fill=mbkasten,draw=black] coordinates {(0,52.2) (1,39.6) (2,35.3) (3,19.2) (4,49.3) (5,85.1) (6,36.1) (7,82.4) (8,21.7) (9,54.6) (10,64.1) (11,42.1)};\end{axis}\end{tikzpicture}\par
\rechenplatz{1}
\fusshilfe{\mbox{12:~65,9 mm}}
\end{pfaufg}
\begin{pfaufg}{P10 ’23}{13.}{}
Auf der Welt werden wohl über 7 000 Sprachen gesprochen. Die Tabelle nennt die fünf Sprachen mit den meisten Muttersprachlern (Anzahl in Millionen): Arabisch 290, Chinesisch 1 300, Englisch 500, Hindi 525, Spanisch 389. Gib für diese fünf Werte das Minimum, das Maximum und die Spannweite an.
\par\smallskip\noindent \sachtabelle{lr}{Sprache & Muttersprachler in Millionen}{Arabisch & 290\\ Chinesisch & 1 300\\ Englisch & 500\\ Hindi & 525\\ Spanisch & 389}\par
\fusshilfe{\mbox{13:~Minimum 290 Mio.; Maximum 1300 Mio.; Spannweite 1010 Mio.}}
\end{pfaufg}
\pfgruppe{}
\begin{pfaufg}{P10 ’20}{14.}{}
Das Säulendiagramm zeigt die Besucherzahlen eines Freibads an den Tagen der ersten Ferienwoche. Die beiden folgenden Aussagen sind richtig: „Die Besucherzahlen der ersten Ferienwoche haben eine Spannweite von 150.“ „Am Sonntag gab es ein Drittel mehr Besucher als am Mittwoch.“ Weise beide Aussagen rechnerisch nach.
\par\smallskip\noindent \begin{tikzpicture}\begin{axis}[ybar,width=11.1cm,height=4.6cm,ymin=0,ymax=336.00,symbolic x coords={Mo,Di,Mi,Do,Fr,Sa,So},xtick=data,enlarge x limits=0.08,bar width=6mm,ylabel={Besucher je Wochentag},ylabel style={font=\small},tick label style={font=\small},nodes near coords,every node near coord/.append style={font=\scriptsize,/pgf/number format/.cd,use comma,1000 sep={}},ymajorgrids,grid style={mbgitter}]\addplot[fill=mbkasten,draw=black] coordinates {(Mo,170.0) (Di,130.0) (Mi,210.0) (Do,180.0) (Fr,190.0) (Sa,240.0) (So,280.0)};\end{axis}\end{tikzpicture}\par
\rechenplatz{3}
\fusshilfe{\mbox{14:~280 $-$ 130 = 150; 210 + 210 : 3 = 280}}
\end{pfaufg}
\pfgruppe{}
\begin{pfaufg}{HH ’26}{15.}{}
Kurswahl von 93 Jugendlichen: Hörbücher 13, Comics 10, Fußball 22, Spiele 16, Angeln 12, Ballspiele 20. Gib den beliebtesten Kurs an, berechne die Spannweite und den Durchschnitt pro Kurs.
\par\smallskip\noindent \sachtabelle{lr}{Kurs & Anzahl}{}\par
\fusshilfe{\mbox{15:~Fußball; 12; 15,5}}
\end{pfaufg}
\pfgruppe{}
\begin{pfaufg}{P10 ’21}{16.}{}
Das Säulendiagramm gibt an, wie viele Container (in Millionen) in den Jahren 2009 bis 2017 im Hamburger Hafen be- und entladen wurden. Berechne die Spannweite und den Durchschnitt (das arithmetische Mittel) dieser Containerzahlen. Berechne, um wie viel Prozent die Containerzahl von 2010 auf 2011 gestiegen ist.
\par\smallskip\noindent \begin{tikzpicture}\begin{axis}[ybar,width=13.7cm,height=4.6cm,ymin=0,ymax=11.64,symbolic x coords={2009,2010,2011,2012,2013,2014,2015,2016,2017},xtick=data,enlarge x limits=0.08,bar width=6mm,ylabel={Container in Millionen},ylabel style={font=\small},tick label style={font=\small},nodes near coords,every node near coord/.append style={font=\scriptsize,/pgf/number format/.cd,use comma,1000 sep={}},ymajorgrids,grid style={mbgitter}]\addplot[fill=mbkasten,draw=black] coordinates {(2009,7.0) (2010,7.9) (2011,9.0) (2012,8.9) (2013,9.3) (2014,9.7) (2015,8.8) (2016,8.9) (2017,8.8)};\end{axis}\end{tikzpicture}\par
\rechenplatz{4}
\fusshilfe{\mbox{16:~2,7 Mio.; 8,7 Mio.; $\tfrac{11}{79}$ $\approx$ 13,9 \%}}
\end{pfaufg}
\pfsteckt{steckt auch in Nr. 45 (P10 ’25)}
\pfstufekopf[smittelwert]{Mittelwert}{in 4 der letzten 5 Prüfungen}
\pfgruppe{}
\begin{pfaufg}{P10 ’16}{17.}{}
Gib das arithmetische Mittel (den Durchschnitt) der Werte 8; 40; 60 an.
\rechenplatz{1}
\fusshilfe{\mbox{17:~36}}
\end{pfaufg}
\pfgruppe{}
\begin{pfaufg}{P10 ’20}{18.}{}
Ein Schwimmbad zählt in der ersten Woche der Sommerferien jeden Tag seine Besucher. Am Donnerstag kamen 180 Besucher. Im Balkendiagramm fehlt der Balken für Donnerstag noch. Berechne, wie viele Besucher in dieser Woche durchschnittlich an einem Tag im Schwimmbad waren.
\par\smallskip\noindent \begin{tikzpicture}\begin{axis}[xmin=0,xmax=400,xtick distance=100,tick label style={font=\small},xticklabel style={/pgf/number format/use comma,/pgf/number format/1000 sep={\,},/pgf/number format/fixed,/pgf/number format/precision=0},xmajorgrids,grid style={mbgitter},ytick={0,1,2,3,4,5,6},yticklabels={{Mo},{Di},{Mi},{Do},{Fr},{Sa},{So}},ymin=-0.5,ymax=6.5,minor x tick num=1,xminorgrids,nodes near coords,every node near coord/.append style={font=\scriptsize,/pgf/number format/use comma,/pgf/number format/1000 sep={\,},/pgf/number format/fixed,/pgf/number format/precision=0},bar width=4mm,xtick pos=bottom,ytick pos=left,xbar,y dir=reverse,width=11cm,height=6.8cm,xlabel={Anzahl der Besucher}]\addplot[fill=mbkasten,draw=black] coordinates {(170,0) (130,1) (210,2) (190,4) (240,5) (280,6)};\end{axis}\end{tikzpicture}\par
\rechenplatz{2}
\fusshilfe{\mbox{18:~200 Besucher}}
\end{pfaufg}
\pfgruppe{}
\begin{pfaufg}{P10 ’17}{19.}{}
Gib den Durchschnitt (das arithmetische Mittel) dieser Längen an: 1,8 m; 1,7 m; 1,6 m; 1,7 m
\rechenplatz{1}
\fusshilfe{\mbox{19:~1,7 m}}
\end{pfaufg}
\begin{pfaufg}{P10 ’21}{20.}{}
Jonas springt im Weitsprung 4,08 m, 3,88 m, 3,92 m und 4,12 m. Gib seine durchschnittliche Weite an.
\rechenplatz{2}
\fusshilfe{\mbox{20:~4,00 m}}
\end{pfaufg}
\pfgruppe{}
\begin{pfaufg}{P10 ’15}{21.}{}
Eine Fußballmannschaft hatte in einer Saison 17 Heimspiele. Zu diesen Spielen kamen zusammen 680 353 Zuschauer. Berechne, wie viele Zuschauer im Durchschnitt bei einem Heimspiel waren. Runde sinnvoll.
\rechenplatz{1}
\fusshilfe{\mbox{21:~680 $\tfrac{353}{17}$ $\approx$ 40 000 Zuschauer (40 021)}}
\end{pfaufg}
\pfgruppe{}
\begin{pfaufg}{NI ’24}{22.}{}
Die Tabelle nennt die Einwohner der neun Nachbarländer Deutschlands in Millionen. Berechne den Durchschnitt.
\par\smallskip\noindent \sachtabelle{l}{Dänemark 5,9}{Niederlande 17,6\\ Belgien 11,7\\ Luxemburg 0,65\\ Frankreich 65\\ Schweiz 8,9\\ Österreich 9\\ Tschechien 10,7\\ Polen 41 (Einwohner in Mio.)}\par
\fusshilfe{\mbox{22:~170,45 Mio. : 9 $\approx$ 18,94 Mio.}}
\end{pfaufg}
\pfgruppe{}
\begin{pfaufg}{P10 ’24}{23.}{}
Das Säulendiagramm zeigt, wie viel Niederschlag 2021 in Brandenburg in jedem Monat gefallen ist (in mm). Im ganzen Jahr 2021 fielen 581,7 mm Niederschlag. Zeige, dass im Durchschnitt pro Monat etwa 48,5 mm fielen. Ermittle, um wie viel Prozent der Niederschlag im April unter diesem Monatsdurchschnitt lag.
\par\smallskip\noindent \begin{tikzpicture}\begin{axis}[ymin=0,ymax=100,ytick distance=20,tick label style={font=\small},xtick pos=bottom,ytick pos=left,yticklabel style={/pgf/number format/use comma,/pgf/number format/1000 sep={\,},/pgf/number format/fixed,/pgf/number format/precision=1},ymajorgrids,grid style={mbgitter},xtick={0,1,2,3,4,5,6,7,8,9,10,11},xticklabels={{Jan},{Feb},{Mär},{Apr},{Mai},{Jun},{Jul},{Aug},{Sep},{Okt},{Nov},{Dez}},xmin=-0.5,xmax=11.5,bar width=5mm,nodes near coords,every node near coord/.append style={font=\scriptsize,/pgf/number format/use comma,/pgf/number format/1000 sep={\,},/pgf/number format/fixed,/pgf/number format/precision=1},ybar,width=13.5cm,height=5.2cm,ylabel={Niederschlag in mm},ylabel style={font=\small}]\addplot[fill=mbkasten,draw=black] coordinates {(0,52.2) (1,39.6) (2,35.3) (3,19.2) (4,49.3) (5,85.1) (6,36.1) (7,82.4) (8,21.7) (9,54.6) (10,64.1) (11,42.1)};\end{axis}\end{tikzpicture}\par
\rechenplatz{3}
\fusshilfe{\mbox{23:~$\approx$ 48,5 mm; $\approx$ 60,4 \% unter dem Durchschnitt}}
\end{pfaufg}
\pfsteckt{steckt auch in Nr. 16 (P10 ’21); Nr. 10 (P10 ’22); Nr. 46 (P10 ’25)}
\pfstufekopf[swinkelberechnenundbe]{Winkel berechnen und beschriften}{in 3 der letzten 5 Prüfungen}
\pfgruppe{}
\begin{pfaufg}{}{24.}{}
Ein Satz lautet: „Ein Viertel der Befragten fährt Rad, das sind $25\,\%$.“ Kreuze an, was die Zahl $25$ hier ist.
\pfkreuzzeile{Anteil in \%}\pfkreuzzeile{Winkel in Grad}
\fusshilfe{\mbox{24:~Anteil in \%}}
\end{pfaufg}
\begin{pfaufg}{}{25.}{}
Ein Satz lautet: „$40$ von $100$ Kindern nennen Pizza, das sind $40\,\%$.“ Kreuze an, was die zweite Zahl $40$ hier ist.
\pfkreuzzeile{Anteil in \%}\pfkreuzzeile{Winkel in Grad}
\fusshilfe{\mbox{25:~Anteil in \%}}
\end{pfaufg}
\pfgruppe{}
\begin{pfaufg}{P10 ’21}{26.}{}
Die Tabelle gibt an, welche Güter im Hamburger Hafen be- und entladen werden (Anteile in \%): Schrott 17, Flüssigkeiten ?, Container 66, Getreide 6, Sonstiges 1, gesamt 100. Ergänze den fehlenden Anteil für Flüssigkeiten. Schreib an jeden Sektor des Kreisdiagramms das passende Gut.
\par\smallskip\noindent \begin{tikzpicture}[line width=0.6pt,baseline=(current bounding box.north)]\begin{scope}\clip (0,0) -- (68.4:2.0) arc (68.4:90:2.0) -- cycle;\foreach \i in {-14,...,14}{\foreach \j in {-14,...,14}{\fill (\i*0.16,\j*0.16) circle (0.5pt);}}\end{scope}\draw (0,0) -- (90:2.0);\begin{scope}\clip (0,0) -- (32.4:2.0) arc (32.4:68.4:2.0) -- cycle;\foreach \i in {-12,...,12}{\draw[line width=0.3pt] (\i*0.18-2.5,-2.5) -- ++(5,5);}\end{scope}\draw (0,0) -- (68.4:2.0);\fill[black!60] (0,0) -- (-28.8:2.0) arc (-28.8:32.4:2.0) -- cycle;\draw (0,0) -- (32.4:2.0);\draw (0,0) -- (-28.8:2.0);\fill[black!10] (0,0) -- (-270:2.0) arc (-270:-32.4:2.0) -- cycle;\draw (0,0) -- (-32.4:2.0);\draw (0,0) circle (2.0); \fill (0,0) circle (1.3pt);\draw[line width=0.4pt] (79.2:1.40) -- (79.2:2.40);\draw[line width=0.4pt] ($(79.2:2.40)+(0,-0.00)$) -- ++(1.8,0);\draw[line width=0.4pt] (50.4:1.40) -- (50.4:2.40);\draw[line width=0.4pt] ($(50.4:2.40)+(0,-0.00)$) -- ++(1.8,0);\draw[line width=0.4pt] (1.8:1.40) -- (1.8:2.40);\draw[line width=0.4pt] ($(1.8:2.40)+(0,-0.00)$) -- ++(1.8,0);\draw[line width=0.4pt] (-30.6:1.40) -- (-30.6:2.40);\draw[line width=0.4pt] ($(-30.6:2.40)+(0,-0.00)$) -- ++(1.8,0);\draw[line width=0.4pt] (-151.2:1.40) -- (-151.2:2.40);\draw[line width=0.4pt] ($(-151.2:2.40)+(0,-0.00)$) -- ++(-1.8,0);\end{tikzpicture}\par
\rechenplatz{1}
\fusshilfe{\mbox{26:~10 \%}}
\end{pfaufg}
\pfgruppe{}
\begin{pfaufg}{P10 ’22}{27.}{}
750 Jugendliche wurden gefragt, wie oft sie Comics lesen. Ergebnis: mehrmals pro Woche 39 \%, einmal pro Woche 18 \%, einmal im Monat 27 \%, nie 16 \%. Im Kreisdiagramm ist nur der Sektor „nie“ beschriftet. Beschrifte die drei anderen Sektoren. Berechne den Mittelpunktswinkel für „nie“.
\par\smallskip\noindent \begin{tikzpicture}[line width=0.6pt,baseline=(current bounding box.north)]\draw (0,0) -- (90:2.0);\node[font=\small,anchor=241.2] at (61.2:2.15) {nie};\draw (0,0) -- (32.4:2.0);\draw (0,0) -- (-32.4:2.0);\draw (0,0) -- (-129.6:2.0);\draw (0,0) circle (2.0); \fill (0,0) circle (1.3pt);\draw[line width=0.4pt] (0:1.40) -- (0:2.40);\draw[line width=0.4pt] ($(0:2.40)+(0,-0.00)$) -- ++(1.8,0);\draw[line width=0.4pt] (-81:1.40) -- (-81:2.40);\draw[line width=0.4pt] ($(-81:2.40)+(0,-0.00)$) -- ++(1.8,0);\draw[line width=0.4pt] (-199.8:1.40) -- (-199.8:2.40);\draw[line width=0.4pt] ($(-199.8:2.40)+(0,-0.00)$) -- ++(-1.8,0);\end{tikzpicture}\par
\rechenplatz{2}
\fusshilfe{\mbox{27:~57,6°}}
\end{pfaufg}
\pfgruppe{}
\begin{pfaufg}{P10 ’24}{\llap{\textasteriskcentered\,}28.}{}
Familie Becker verbraucht am Tag etwa 508 Liter Leitungswasser: 1 Trinken und Kochen 12 l, 2 Körperpflege 256 l, 3 Toilettenspülung 136 l, 4 Wäsche waschen 60 l, 5 Gartenpflege 16 l, 6 Sonstiges 28 l. Das Kreisdiagramm zeigt diese Anteile; beschriftet sind nur „1 Trinken und Kochen“ und „5 Gartenpflege“. Beschrifte den Sektor für „Wäsche waschen“. Berechne den Mittelpunktswinkel für den Anteil „Trinken und Kochen“ am gesamten Verbrauch.
\par\smallskip\noindent \begin{tikzpicture}[line width=0.6pt,baseline=(current bounding box.north)]\draw (0,0) -- (90:2.0);\draw (0,0) -- (-91.4:2.0);\draw (0,0) -- (-111.2:2.0);\fill[black!25] (0,0) -- (-165:2.0) arc (-165:-153.7:2.0) -- cycle;\draw (0,0) -- (-153.7:2.0);\node[font=\small,anchor=20.65] at (-159.35:2.15) {5 Gartenpflege};\draw (0,0) -- (-165:2.0);\fill[black!25] (0,0) -- (-269.9:2.0) arc (-269.9:-261.4:2.0) -- cycle;\draw (0,0) -- (-261.4:2.0);\node[font=\small,anchor=274.35] at (-265.65:2.15) {1 Trinken und Kochen};\draw (0,0) circle (2.0); \fill (0,0) circle (1.3pt);\draw[line width=0.4pt] (-0.7:1.40) -- (-0.7:2.40);\draw[line width=0.4pt] ($(-0.7:2.40)+(0,-0.00)$) -- ++(1.8,0);\draw[line width=0.4pt] (-101.3:1.40) -- (-101.3:2.40);\draw[line width=0.4pt] ($(-101.3:2.40)+(0,-0.00)$) -- ++(-1.8,0);\draw[line width=0.4pt] (-132.45:1.40) -- (-132.45:2.40);\draw[line width=0.4pt] ($(-132.45:2.40)+(0,-0.00)$) -- ++(-1.8,0);\draw[line width=0.4pt] (-213.2:1.40) -- (-213.2:2.40);\draw[line width=0.4pt] ($(-213.2:2.40)+(0,-0.00)$) -- ++(-1.8,0);\end{tikzpicture}\par
\rechenplatz{2}
\fusshilfe{\mbox{28:~$\approx$ 43°) beschriften; $\approx$ 8,5°}}
\end{pfaufg}
\pfgruppe{}
\begin{pfaufg}{P10 ’17}{29.}{}
Die Tabelle zeigt, als wievieltes Kind ihrer Mutter die Kinder des Jahrgangs 2013 in Deutschland zur Welt kamen (Anteil in Prozent): erstes Kind 49,4 \%, zweites Kind 34,4 \%, drittes Kind 11,2 \%, viertes oder weiteres Kind 5,0 \%. Im Kreisdiagramm ist einer dieser Anteile grau gefärbt. Schreib an den grauen Sektor, zu welchem Anteil er gehört. Berechne, wie groß der Mittelpunktswinkel für „drittes Kind“ sein muss. Zeichne diesen Sektor in den Kreis ein und beschrifte ihn.
\par\smallskip\noindent \begin{tikzpicture}[line width=0.6pt,baseline=(current bounding box.north)]\fill[black!25] (0,0) -- (-180:2.0) arc (-180:-162:2.0) -- cycle;\draw (0,0) -- (-180:2.0);\draw (0,0) circle (2.0); \fill (0,0) circle (1.3pt);\draw[line width=0.4pt] (-171:1.40) -- (-171:2.40);\draw[line width=0.4pt] ($(-171:2.40)+(0,-0.00)$) -- ++(-1.8,0);\draw[line width=0.4pt] ($(-171:2.40)+(0,-0.55)$) -- ++(-1.8,0);\draw[line width=0.4pt] ($(-171:2.40)+(0,-1.10)$) -- ++(-1.8,0);\end{tikzpicture}\par
\rechenplatz{2}
\fusshilfe{\mbox{29:~$\approx$ 40°; Sektor 40° „drittes Kind“}}
\end{pfaufg}
\pfsteckt{steckt auch in Nr. 49 (P10 ’25)}
\pfstufekopf[sergnzen]{ergänzen}{in 2 der letzten 5 Prüfungen}
\pfgruppe{}
\begin{pfaufg}{}{30.}{}
Die y-Achse (senkrechte Achse) eines Säulendiagramms ist mit $0$, $20$, $40$, $60$, $80$ beschriftet. Kreuze an, ob die Achse bei null beginnt. Wenn nicht, trage den Startwert ein.
\pfkreuzzeile{ja, bei null}\pfkreuzzeile{nein, Startwert}
\fusshilfe{\mbox{30:~ja, bei null}}
\end{pfaufg}
\begin{pfaufg}{}{31.}{}
Sieh dir die y-Achse des Säulendiagramms an. Wie viel ist eine Hilfslinie wert? Gib nur diesen Wert an.
\par\smallskip\noindent \saeulenab[ymax=2,ystep=0.5,yfein=0.1,ylabel=Länge in m]{A/1.3,B/0.7,C/1.8}\par
\par\noindent \leerfeld[je]  Linie\par
\fusshilfe{\mbox{31:~$0{,}1$ je Linie}}
\end{pfaufg}
\pfgruppe{}
\begin{pfaufg}{P10 ’20}{32.}{}
Ein Schwimmbad zählt in der ersten Woche der Sommerferien jeden Tag seine Besucher. Am Donnerstag kamen 180 Besucher. Im Balkendiagramm fehlt der Balken für Donnerstag noch. Zeichne im Diagramm den Balken für Donnerstag ein.
\par\smallskip\noindent \begin{tikzpicture}\begin{axis}[xmin=0,xmax=400,xtick distance=100,tick label style={font=\small},xticklabel style={/pgf/number format/use comma,/pgf/number format/1000 sep={\,},/pgf/number format/fixed,/pgf/number format/precision=0},xmajorgrids,grid style={mbgitter},ytick={0,1,2,3,4,5,6},yticklabels={{Mo},{Di},{Mi},{Do},{Fr},{Sa},{So}},ymin=-0.5,ymax=6.5,minor x tick num=1,xminorgrids,nodes near coords,every node near coord/.append style={font=\scriptsize,/pgf/number format/use comma,/pgf/number format/1000 sep={\,},/pgf/number format/fixed,/pgf/number format/precision=0},bar width=4mm,xtick pos=bottom,ytick pos=left,xbar,y dir=reverse,width=11cm,height=6.8cm,xlabel={Anzahl der Besucher}]\addplot[fill=mbkasten,draw=black] coordinates {(170,0) (130,1) (210,2) (190,4) (240,5) (280,6)};\end{axis}\end{tikzpicture}\par
\rechenplatz{1}
\fusshilfe{\mbox{32:~Balken bei Do bis 180}}
\end{pfaufg}
\begin{pfaufg}{P10 ’18}{33.}{}
1000 Jugendliche haben ihr Lieblingsessen genannt: Schnitzel 80, Pizza 90, Kartoffelgerichte 160, Gemüsegerichte 160, Fischgerichte 110; die Zahlen für Pasta und Salat fehlen noch. Die fünf bekannten Zahlen sind in einem Balkendiagramm dargestellt. Die Achse „Anzahl“ hat aber noch keine Einteilung, und die Balken sind nur mit A bis E bezeichnet. Trag an der Achse eine passende Einteilung ein. Schreib zu jedem Balken das passende Gericht.
\par\smallskip\noindent \begin{tikzpicture}[x=0.45cm,y=0.45cm]\draw[mbgitter,line width=0.3pt,step=1] (0,0) grid (18,11);\draw[line width=0.7pt] (0,0) -- (0,11);\draw[line width=0.7pt,->] (0,0) -- (11,0);\fill[mbkasten,draw=black] (0,8.60) rectangle (5.5,9.40);\node[left,font=\small] at (0,9) {A};\draw[line width=0.4pt] (11,8.6) -- (17,8.6);\fill[mbkasten,draw=black] (0,6.60) rectangle (4.0,7.40);\node[left,font=\small] at (0,7) {B};\draw[line width=0.4pt] (11,6.6) -- (17,6.6);\fill[mbkasten,draw=black] (0,4.60) rectangle (4.5,5.40);\node[left,font=\small] at (0,5) {C};\draw[line width=0.4pt] (11,4.6) -- (17,4.6);\fill[mbkasten,draw=black] (0,2.60) rectangle (8.0,3.40);\node[left,font=\small] at (0,3) {D};\draw[line width=0.4pt] (11,2.6) -- (17,2.6);\fill[mbkasten,draw=black] (0,0.60) rectangle (8.0,1.40);\node[left,font=\small] at (0,1) {E};\draw[line width=0.4pt] (11,0.6) -- (17,0.6);\end{tikzpicture}\par
\rechenplatz{2}
\fusshilfe{\mbox{33:~20 je Gitterlinie}}
\end{pfaufg}
\pfgruppe{}
\begin{pfaufg}{P10 ’23}{34.}{}
In einem Säulendiagramm sind zwei Werte aus der Tabelle schon als Säulen gezeichnet; die erste Säule gehört zu Spanisch. Trag an der y-Achse eine passende Einteilung ein. Schreib unter die zweite Säule, zu welcher Sprache sie gehört. Zeichne die Säule für Arabisch ein.
\par\smallskip\noindent \begin{tikzpicture}[x=0.45cm,y=0.45cm]\draw[mbgitter,line width=0.3pt] (0,0) grid (16,14);\draw[line width=0.7pt,->] (0,0) -- (0,14.6) node[above,font=\small] {Anzahl in Mio.};\draw[line width=0.7pt,->] (0,0) -- (16.6,0) node[right,font=\small] {Sprache};\fill[mbkasten,draw=black] (2,0) rectangle (4,7.8);\node[below=2pt,font=\small] at (3,0) {Spanisch};\fill[mbkasten,draw=black] (7,0) rectangle (9,10.5);\node[below=2pt,font=\small] at (8,0) {\leerzelle};\node[below=2pt,font=\small] at (13,0) {Arabisch};\end{tikzpicture}\par
\fusshilfe{\mbox{34:~50 Mio.; Hindi; 5,8 Kästchen hoch}}
\end{pfaufg}
\pfstufekopf[saussageprfen]{Aussage prüfen}{in 1 der letzten 5 Prüfungen}
\pfgruppe{}
\begin{pfaufg}{}{35.}{}
Die y-Achse eines Säulendiagramms ist mit $0$, $50$, $100$, $150$ beschriftet. Kreuze an, wo sie beginnt.
\pfkreuzzeile{bei null}\pfkreuzzeile{nicht bei null}
\fusshilfe{\mbox{35:~bei null}}
\end{pfaufg}
\begin{pfaufg}{}{36.}{}
Sieh dir die y-Achse des Liniendiagramms an. Kreuze an, wo sie beginnt. Wenn sie nicht bei null beginnt, schreibe den Startwert auf.
\pfkreuzzeile{bei null}\pfkreuzzeile{nicht bei null}
\par\smallskip\noindent \begin{tikzpicture}\begin{axis}[axis lines=left,width=9cm,height=6cm,xmin=2019,xmax=2025,ymin=80,ymax=96,xtick distance=1,ytick distance=2,grid=both,grid style={mbgitter},tick label style={font=\scriptsize},/pgf/number format/1000 sep={},xlabel={Jahr},ylabel={Mitglieder},label style={font=\small}]\addplot[only marks,mark=*,mark size=1.5pt] coordinates {(2020,84) (2021,86) (2022,90) (2023,91) (2024,94)};\end{axis}\end{tikzpicture}\par
\fusshilfe{\mbox{36:~nicht bei null, Startwert $80$}}
\end{pfaufg}
\pfgruppe{}
\begin{pfaufg}{P10 ’19}{\llap{\textasteriskcentered\,}37.}{}
2015 erschien ein Zeitungsartikel mit der Überschrift „Bücherlesen bald out?“. Darin heißt es: Die Buchhändler klagen, dass jedes Jahr dramatisch weniger Jugendliche Bücher lesen. Man fürchte, dass schon bald kein einziger Jugendlicher mehr zu einem Buch greift. Neben dem Text steht das Säulendiagramm. Schüler einer 10. Klasse halten zwei Aussagen des Artikels für falsch. Nenne eine davon und berichtige sie.
\par\smallskip\noindent \begin{tikzpicture}\begin{axis}[ymin=30,ymax=50,ytick distance=5,tick label style={font=\small},xtick pos=bottom,ytick pos=left,yticklabel style={/pgf/number format/use comma,/pgf/number format/1000 sep={\,},/pgf/number format/fixed,/pgf/number format/precision=0},ymajorgrids,grid style={mbgitter},xtick={0,1,2,3,4},xticklabels={{2011},{2012},{2013},{2014},{2015}},xmin=-0.5,xmax=4.5,bar width=7mm,nodes near coords,every node near coord/.append style={font=\scriptsize,/pgf/number format/use comma,/pgf/number format/1000 sep={\,},/pgf/number format/fixed,/pgf/number format/precision=0},ybar,width=8.0cm,height=5.2cm,ylabel={Bücher lesende Jugendliche in \%},ylabel style={font=\small}]\addplot[fill=mbkasten,draw=black] coordinates {(0,44) (1,42) (2,40) (3,39) (4,36)};\end{axis}\end{tikzpicture}\par
\rechenplatz{1}
\fusshilfe{\mbox{37:~„extrem abnehmen“ falsch}}
\end{pfaufg}
\pfgruppe{}
\begin{pfaufg}{}{38.}{}
Eine Stadt zählt jedes Jahr die Fahrraddiebstähle: 2020 $410$, 2021 $380$, 2022 $395$, 2023 $350$. Die Zeitung schreibt: „Jedes Jahr gibt es weniger Diebstähle. 2030 gibt es keine mehr.“ Prüfe beide Teile der Aussage.
\rechenplatz{4}
\end{pfaufg}
\pfgruppe{}
\begin{pfaufg}{P10 ’22}{39.}{}
Von 2017 bis 2021 wurden jedes Jahr 1 200 Jugendliche gefragt, wie oft sie Bücher lesen. Das Säulendiagramm zeigt, wie viele von ihnen mehrmals pro Woche ein Buch lesen. Ben schaut auf das Diagramm und sagt: „Jedes Jahr lesen weniger Jugendliche Bücher. Von 2018 bis 2021 ist ihre Zahl um mehr als die Hälfte gesunken.“ Entscheide, ob Ben recht hat. Begründe deine Entscheidung.
\par\smallskip\noindent \begin{tikzpicture}\begin{axis}[ybar,width=8.5cm,height=4.6cm,ymin=390,ymax=508.75,symbolic x coords={2017,2018,2019,2020,2021},xtick=data,enlarge x limits=0.08,bar width=6mm,ylabel={Anzahl},ylabel style={font=\small},tick label style={font=\small},nodes near coords,every node near coord/.append style={font=\scriptsize,/pgf/number format/.cd,use comma,1000 sep={}},ymajorgrids,grid style={mbgitter}]\addplot[fill=mbkasten,draw=black] coordinates {(2017,468.0) (2018,432.0) (2019,485.0) (2020,470.0) (2021,400.0)};\end{axis}\end{tikzpicture}\par
\rechenplatz{1}
\fusshilfe{\mbox{39:~2019 Anstieg 432 auf 485; $\approx$ 7,4 \%), nicht um die Hälfte}}
\end{pfaufg}
\pfgruppe{}
\begin{pfaufg}{P10 ’18}{\llap{\textasteriskcentered\,}40.}{}
Frauen wurden gefragt, wie oft sie Fleisch oder Wurst essen. Das Säulendiagramm zeigt die Anteile der Antworten. Prüfe die beiden Aussagen zum Säulendiagramm. Kreuze jeweils an, ob sie richtig, falsch oder nicht entscheidbar ist, und begründe. Aussage 1:
\par\smallskip\noindent \begingroup\renewcommand{\arraystretch}{1.5}\begin{tabular}{|p{5.5cm}|c|c|c|}\hline Aussage & richtig & falsch & nicht entscheidbar \\ \hline „Nicht einmal die Hälfte der Frauen isst 4- bis 6-mal pro Woche Fleisch oder Wurst.“ & $\square$ & $\square$ & $\square$ \\ \hline \multicolumn{4}{|l|}{\rule{0pt}{3.2ex}\footnotesize\color{mbgrau}Begründung:} \\ \hline „1- bis 3-mal pro Woche Fleisch oder Wurst isst jede 15. Frau.“ & $\square$ & $\square$ & $\square$ \\ \hline \multicolumn{4}{|l|}{\rule{0pt}{3.2ex}\footnotesize\color{mbgrau}Begründung:} \\ \hline\end{tabular}\endgroup\par
\fusshilfe{\mbox{40:~Aussage 1 falsch (55 \% > 50 \%); Aussage 2 falsch}}
\end{pfaufg}
\pfstufekopf[smedian]{Median}{in 1 der letzten 5 Prüfungen}
\pfgruppe{}
\begin{pfaufg}{P10 ’15}{41.}{}
Gib den Median (Zentralwert) dieser Messwerte an: 5 °C; 7 °C; 3 °C; 8 °C; 1 °C; 8 °C; 5 °C
\rechenplatz{1}
\fusshilfe{\mbox{41:~5 °C}}
\end{pfaufg}
\pfgruppe{}
\begin{pfaufg}{P10 ’24}{42.}{}
Ermittle den Median (Zentralwert) dieser Datenreihe: 20 °C; 17 °C; 21 °C; 18 °C; 21 °C; 11 °C
\rechenplatz{2}
\fusshilfe{\mbox{42:~19 °C}}
\end{pfaufg}

\par\Needspace*{0.55\textheight}\pfstufekopf[serkennen1]{Was ist gesucht?}{}
\begin{pfaufg}{}{43.}{}
Was ist gesucht? Kreuze an. Rechne nicht.
\par\smallskip\noindent \begingroup\renewcommand{\arraystretch}{1.4}\begin{tabular}{|>{\raggedright\arraybackslash}p{6.3cm}|>{\centering\arraybackslash\hyphenpenalty=0}p{1.55cm}|>{\centering\arraybackslash\hyphenpenalty=0}p{1.55cm}|>{\centering\arraybackslash\hyphenpenalty=0}p{1.55cm}|}\hline  & {\scriptsize\hspace{0pt}Mittelwert} & {\scriptsize\hspace{0pt}Minimum} & {\scriptsize\hspace{0pt}Median} \\ \hline Wie viele Besucher kamen im Durchschnitt pro Tag? & $\square$ & $\square$ & $\square$ \\ \hline Welche Sprache hat die meisten, welche die wenigsten Muttersprachler? & $\square$ & $\square$ & $\square$ \\ \hline Temperaturen: 20, 17, 21, 18, 21, 11 °C. Bestimme den Median. & $\square$ & $\square$ & $\square$ \\ \hline Liegt der Benzinpreis von Montag bis Freitag im Mittel unter 1,80 €? & $\square$ & $\square$ & $\square$ \\ \hline Wie groß ist der Unterschied zwischen dem nassesten und dem trockensten Monat 2021? & $\square$ & $\square$ & $\square$ \\ \hline Messwerte: 5, 7, 3, 8, 1, 8, 5 °C. Welcher Wert steht in der Mitte, wenn man sie ordnet? & $\square$ & $\square$ & $\square$ \\ \hline\end{tabular}\endgroup\par
\end{pfaufg}
\pfstufekopf[srckwrtsfehlenderwert]{rückwärts: fehlender Wert}{in 1 der letzten 5 Prüfungen}
\pfgruppe{}
\begin{pfaufg}{P10 ’22}{44.}{}
In einer Woche lag die Temperatur im Durchschnitt bei 20 °C. Ergänze die fehlende Temperatur am Samstag.
\par\smallskip\noindent \sachtabelle{lcccccc}{Mo & Di & Mi & Do & Fr & Sa & So}{21 °C & 20 °C & 19 °C & 22 °C & 20 °C & \leerzelle & 20 °C}\par
\fusshilfe{\mbox{44:~18 °C}}
\end{pfaufg}
\pfstufekopf[saussagenprfenauswirk]{Aussagen prüfen, Auswirkung erklären}{in 1 der letzten 5 Prüfungen}
\pfgruppe{}
\begin{pfaufg}{P10 ’25}{45.}{}
In einer Kita-Gruppe sind die Kinder so alt (in Jahren): 5, 2, 1, 4, 2, 4, 2, 4, 2, 2. Eine der beiden Aussagen ist falsch: Kreuze die falsche Aussage an und berichtige sie.
\par\smallskip\noindent \begingroup\renewcommand{\arraystretch}{1.5}\begin{tabular}{|p{7.8cm}|c|p{3.4cm}|}\hline Aussage & falsch & Korrektur \\ \hline Die Spannweite beträgt 4 Jahre. & $\square$ &  \\ \hline Der Modalwert beträgt 4 Jahre. & $\square$ &  \\ \hline\end{tabular}\endgroup\par
\fusshilfe{\mbox{45:~4 stimmt}}
\end{pfaufg}
\pfgruppe{}
\begin{pfaufg}{P10 ’25}{\llap{\textasteriskcentered\,}46.}{}
In derselben Kita arbeiten 11 Erziehende. Sie sind so alt (in Jahren): 52, 57, 33, 44, 31, 30, 66, 58, 61, 41, 55. Bestimme das Durchschnittsalter der Erziehenden. Die älteste und die jüngste Person verlassen die Kita. Erkläre, wie sich dadurch das Durchschnittsalter verändert.
\rechenplatz{3}
\fusshilfe{\mbox{46:~48 Jahre; 2 · 48}}
\end{pfaufg}
\pfstufekopf[sausprozentdarstellen]{aus Prozent darstellen}{in 1 der letzten 5 Prüfungen}
\pfgruppe{}
\begin{pfaufg}{P10 ’18}{47.}{}
Frauen wurden gefragt, wie oft sie Fleisch oder Wurst essen. Das Säulendiagramm zeigt die Anteile der Antworten. Stelle die Anteile aus dem Säulendiagramm als Streifendiagramm dar. Der vorgezeichnete Streifen ist 10 cm lang.
\par\smallskip\noindent \begin{tikzpicture}\draw[mbgitter,line width=0.3pt,step=0.5] (-0.5,-0.5) grid (10.5,2);\draw[line width=0.7pt] (0,0) rectangle (10,1.5);\end{tikzpicture}\par
\rechenplatz{2}
\end{pfaufg}
\pfgruppe{}
\begin{pfaufg}{NI ’24}{48.}{}
In Deutschland haben 41 \% Blutgruppe 0, 43 \% Blutgruppe A und 5 \% Blutgruppe AB. Der Rest hat Blutgruppe B. Ergänze den Anteil. Zeichne ein Streifendiagramm (10 cm = 100 \%).
\par\smallskip\noindent \sachtabelle{l}{Blutgruppe 0, A, B, AB}{Anteil 41 \%, 43 \%, leer, 5 \%}
\par\smallskip\noindent \begin{tikzpicture}\draw[mbgitter,line width=0.3pt,step=0.5] (-0.5,-0.5) grid (10.5,2);\draw[line width=0.7pt] (0,0) rectangle (10,1.5);\end{tikzpicture}\par
\fusshilfe{\mbox{48:~B: 11 \%}}
\end{pfaufg}
\pfgruppe{}
\begin{pfaufg}{P10 ’25}{49.}{}
In derselben Kita arbeiten 11 Erziehende. Sie sind so alt (in Jahren): 52, 57, 33, 44, 31, 30, 66, 58, 61, 41, 55. Ermittle, wie viel Prozent der Erziehenden jünger als 50 Jahre sind. Stelle diesen Anteil in einem Kreisdiagramm dar.
\par\smallskip\noindent \kreisleer[2.5]\par
\rechenplatz{3}
\fusshilfe{\mbox{49:~$\approx$ 45,5 \%; $\approx$ 164°}}
\end{pfaufg}
\pfstufekopf[sfalscheneindruckerkl]{falschen Eindruck erklären}{in 1 der letzten 5 Prüfungen}
\pfgruppe{}
\begin{pfaufg}{}{50.}{}
Ein Diagramm zeigt die Werte $1\,020$, $1\,040$ und $1\,060$. Seine Achse beginnt bei $1\,000$. Skizziere das Diagramm mit einer Achse ab $0$. Vergleiche, wie die beiden Diagramme wirken.
\par\smallskip\noindent \saeulen{2021/,2022/,2023/}{1200}{200}{Anzahl}\par
\rechenplatz{2}
\end{pfaufg}
\pfgruppe{}
\begin{pfaufg}{P10 ’14}{\llap{\textasteriskcentered\,}51.}{}
Eine andere Bank zahlt nur 0,25 \% Zinsen im Jahr. Für ihre Werbung „Tolle Zinsen!“ zeichnet sie ein Säulendiagramm, in dem das Guthaben kräftig zu wachsen scheint. Erkläre, mit welchem Trick die Grafik die Entwicklung des Guthabens besser aussehen lässt, als sie ist.
\par\smallskip\noindent \begin{tikzpicture}\begin{axis}[ymin=1060,ymax=1072,ytick distance=2,tick label style={font=\small},xtick pos=bottom,ytick pos=left,yticklabel style={/pgf/number format/use comma,/pgf/number format/1000 sep={\,},/pgf/number format/fixed,/pgf/number format/precision=1},ymajorgrids,grid style={mbgitter},xtick={0,1,2},xticklabels={{nach einem Jahr},{nach zwei Jahren},{nach drei Jahren}},xmin=-0.5,xmax=2.5,bar width=7mm,nodes near coords,every node near coord/.append style={font=\scriptsize,/pgf/number format/use comma,/pgf/number format/1000 sep={\,},/pgf/number format/fixed,/pgf/number format/precision=1},ybar,width=8.2cm,height=5.2cm,ylabel={Guthaben in Euro},ylabel style={font=\small},xticklabel style={text width=2cm,align=center,font=\scriptsize}]\addplot[fill=mbkasten,draw=black] coordinates {(0,1064) (1,1067) (2,1069.5)};\end{axis}\end{tikzpicture}\par
\rechenplatz{1}
\fusshilfe{\mbox{51:~y-Achse beginnt bei 1060 € statt bei 0}}
\end{pfaufg}
\pfsteckt{steckt auch in Nr. 37 (P10 ’19)}
\pfpruefstein{P10 ’26}
\begin{pfaufg}{}{}{}
Eine Tankstelle hat eine Woche lang jeden Tag um 18 Uhr den Benzinpreis notiert.
\par\smallskip\noindent \sachtabelle{lrrrrrrr}{Wochentag & Mo & Di & Mi & Do & Fr & Sa & So}{Benzinpreis in € pro Liter & 1,77 & 1,78 & 1,77 & 1,79 & 1,84 & 1,82 & 1,85}\par
\end{pfaufg}
\begin{pfaufg}{}{a)}{1 BE}
Gib die Spannweite der Benzinpreise in € pro Liter an.
\fusshilfe{\mbox{Prüfstein a):~0,08 € pro Liter}}
\end{pfaufg}
\begin{pfaufg}{}{b)}{2 BE}
Ein Kunde meint: Von Montag bis Freitag hat ein Liter Benzin an dieser Tankstelle im Durchschnitt weniger als 1,80 € gekostet. Zeige mit einer Rechnung, dass er recht hat.
\fusshilfe{\mbox{Prüfstein b):~1,79 € < 1,80 €}}
\end{pfaufg}
\begin{pfaufg}{}{c)}{2 BE}
An einer Tankstelle wurde eine Woche lang täglich um 18 Uhr der Benzinpreis notiert (€ pro Liter): Mo 1,77, Di 1,78, Mi 1,77, Do 1,79, Fr 1,84, Sa 1,82, So 1,85. Berechne die prozentuale Preissteigerung von Donnerstag auf Freitag.
\rechenplatz{3}
\fusshilfe{\mbox{Prüfstein c):~$\tfrac{5}{179}$ $\approx$ 2,8 \%}}
\end{pfaufg}
\begin{pfaufg}{}{d)}{1 BE}
Timo hat die Preise aus der Tabelle in ein Säulendiagramm übertragen. Die Säule für Sonntag fehlt noch. Zeichne sie ein.
\par\smallskip\noindent \begin{tikzpicture}\begin{axis}[ymin=1.75,ymax=1.85,ytick distance=0.02,tick label style={font=\small},xtick pos=bottom,ytick pos=left,yticklabel style={/pgf/number format/use comma,/pgf/number format/1000 sep={\,},/pgf/number format/fixed,/pgf/number format/precision=2},ymajorgrids,grid style={mbgitter},xtick={0,1,2,3,4,5,6},xticklabels={{Mo},{Di},{Mi},{Do},{Fr},{Sa},{So}},xmin=-0.5,xmax=6.5,bar width=7mm,minor y tick num=1,yminorgrids,ybar,width=10.2cm,height=5.2cm,ylabel={Benzinpreis in € pro Liter},ylabel style={font=\small}]\addplot[fill=mbkasten,draw=black] coordinates {(0,1.77) (1,1.78) (2,1.77) (3,1.79) (4,1.84) (5,1.82)};\end{axis}\end{tikzpicture}\par
\rechenplatz{3}
\fusshilfe{\mbox{Prüfstein d):~Säule So bis 1,85 (höchste Säule)}}
\end{pfaufg}
\begin{pfaufg}{}{*e)}{1 BE}
Timo schreibt in einer Bewertung der Tankstelle: „Mein Diagramm beweist: Von Donnerstag auf Freitag ist der Preis auf mehr als das Doppelte gestiegen.“ Begründe, wodurch Timos Diagramm diesen falschen Eindruck erweckt.
\par\smallskip\noindent \sachtabelle{lrrrrrrr}{Wochentag & Mo & Di & Mi & Do & Fr & Sa & So}{Benzinpreis in € pro Liter & 1,77 & 1,78 & 1,77 & 1,79 & 1,84 & 1,82 & 1,85}\par{\footnotesize\color{mbgrau}Säule So: 1,85}\par
\fusshilfe{\mbox{Prüfstein e):~Achse beginnt bei 1,75 € statt bei 0}}
\end{pfaufg}
\end{document}